 
% ============================================================
\chapter{Maximale belading en stabiliteitsgrens}
\label{ch:max-belading}
% ============================================================
 
\section{Inleiding}
\label{sec:max-belading-inleiding}
 
Nu het minimaal benodigde motorvermogen is vastgesteld ($P_{B,\mathrm{inst}} \approx 37$~kW,
zie Sectie~\ref{sec:aandrijflijn}) en de energieopslagsysteemkeuze is afgerond
(Hoofdstuk~\ref{ch:energiekeuze}), kan worden beoordeeld of de gekozen batterijconfiguratie
fysiek past binnen het schip. In dit hoofdstuk wordt bepaald hoeveel extra gewicht aan
\emph{t~Jansje} kan worden toegevoegd voordat de stabiliteitseis wordt overschreden of de
diepgang de rompdiepte bereikt. Twee grenswaarden worden beschouwd:
 
\begin{enumerate}
    \item \textbf{Diepgangsgrens}: de geladen diepgang $T_\mathrm{gel}$ mag de rompdiepte
          $D = 2{,}05$~m niet overschrijden. Met een conservatieve minimale vrijboord van
          $f_\mathrm{min} = 0{,}05$~m geldt:
          \begin{equation}
              T_\mathrm{max} = D - f_\mathrm{min} = 2{,}05 - 0{,}05 = 2{,}00\ \mathrm{m}.
              \label{eq:Tmax}
          \end{equation}
    \item \textbf{Stabiliteitsgrens ($GM$)}: de metacentrische hoogte in geladen conditie
          dient te voldoen aan de IMO-minimumeis van $GM_\mathrm{min} = 0{,}15$~m
          \cite{IMO2008} en aan de aanbevolen ondergrens van $GM_\mathrm{aanbevolen} = 0{,}30$~m
          voor binnenvaartschepen \cite{CCR2015}.
\end{enumerate}
 
\section{Scheepsgeometrische invoer}
\label{sec:max-belading-invoer}
 
\begin{table}[H]
    \centering
    \caption{Invoerparameters maximale beladingsberekening (zoetwater, $\rho = 1\,000\ \mathrm{kg/m^3}$)}
    \label{tab:invoer-maxbelading}
    \begin{tabular}{llrr}
        \toprule
        Parameter & Symbool & Waarde & Eenheid \\
        \midrule
        Lengte overall / waterlijn    & $L$                 & 22{,}22 & m    \\
        Maximale breedte              & $B$                 &  4{,}36 & m    \\
        Rompdiepte (kiel–dek)         & $D$                 &  2{,}05 & m    \\
        Leeggewicht                   & $W_0$               & 73{,}0  & t    \\
        Diepgang leeg                 & $T_0$               &  1{,}50 & m    \\
        Blokcoëfficiënt               & $C_B$               &  0{,}502& –    \\
        Hoofdspantcoëfficiënt         & $C_M$               &  0{,}90 & –    \\
        Waterplancoëfficiënt          & $C_{WP}$            &  0{,}795& –    \\
        Zwaartepunt leeg (geschat)    & $KG_0$              &  1{,}230& m    \\
        Zwaartepunt lading (batterijen)& $z_b$              &  1{,}10 & m    \\
        Minimale vrijboord            & $f_\mathrm{min}$    &  0{,}05 & m    \\
        \bottomrule
    \end{tabular}
\end{table}
 
\section{Analytisch stabiliteitsmodel}
\label{sec:analytisch-model}
 
\subsection{Waterplanoppervlak en tweede moment}
 
Het tweede moment van het waterplanoppervlak om de langscheepse as, benodigd voor $BM$,
wordt berekend via de formule van Schneekluth \cite{Schneekluth1998}:
 
\begin{equation}
    I_T = \frac{L\,B^3}{12}\!\left(0{,}096 + 0{,}89\,C_{WP}^{2}\right)
        = \frac{22{,}22 \times 4{,}36^3}{12}\!\left(0{,}096 + 0{,}89 \times 0{,}795^2\right)
        = 101{,}1\ \mathrm{m}^4.
    \label{eq:IT}
\end{equation}
 
De tonnemaat per meter diepgang (TPM) in zoetwater bedraagt:
 
\begin{equation}
    \mathrm{TPM} = \frac{\rho\,C_{WP}\,L\,B}{1000}
               = \frac{1\,000 \times 0{,}795 \times 22{,}22 \times 4{,}36}{1000}
               = 77{,}0\ \mathrm{t/m}.
    \label{eq:TPM}
\end{equation}
 
\subsection{Betrekkingen als functie van de ladingmassa}
 
Voor een toegevoegde ladingmassa $M_\mathrm{lading}$ (in kilogram) gelden:
 
\begin{align}
    \Delta T &= \frac{M_\mathrm{lading}}{\rho\,C_{WP}\,L\,B}
              = \frac{M_\mathrm{lading}}{77\,004}\ \mathrm{m},
    \qquad
    T_\mathrm{gel} = T_0 + \Delta T,
    \label{eq:deltaT} \\[6pt]
    \nabla_\mathrm{gel} &= \frac{W_0 \times 10^3 + M_\mathrm{lading}}{\rho}
                        = \frac{73\,000 + M_\mathrm{lading}}{1\,000}\ \mathrm{m}^3,
    \label{eq:nabla-gel} \\[6pt]
    KB_\mathrm{gel} &= T_\mathrm{gel}\!\left(0{,}9 - 0{,}3\,C_M - 0{,}1\,C_B\right)
                    = T_\mathrm{gel} \times 0{,}5798\ \mathrm{m},
    \label{eq:KB-gel} \\[6pt]
    BM_\mathrm{gel} &= \frac{I_T}{\nabla_\mathrm{gel}}
                    = \frac{101{,}1}{\nabla_\mathrm{gel}}\ \mathrm{m},
    \label{eq:BM-gel} \\[6pt]
    KG_\mathrm{gel} &= \frac{W_0 \times 10^3 \times KG_0 + M_\mathrm{lading} \times z_b}
                             {W_0 \times 10^3 + M_\mathrm{lading}}
                      = \frac{73\,000 \times 1{,}230 + M_\mathrm{lading} \times 1{,}10}
                             {73\,000 + M_\mathrm{lading}}\ \mathrm{m},
    \label{eq:KG-gel} \\[6pt]
    GM_\mathrm{gel} &= KB_\mathrm{gel} + BM_\mathrm{gel} - KG_\mathrm{gel}.
    \label{eq:GM-gel}
\end{align}
 
\section{Bepaling van de maximale belading}
\label{sec:max-belading-resultaat}
 
\subsection{Diepgangsgrens}
 
De maximaal toelaatbare ladingmassa op basis van $T_\mathrm{max} = 2{,}00$~m:
 
\begin{equation}
    M_{\mathrm{lading,T}} = \rho\,C_{WP}\,L\,B\,(T_\mathrm{max} - T_0)
                          = 77\,004 \times (2{,}00 - 1{,}50)
                          = 38\,502\ \mathrm{kg} \approx \mathbf{38{,}5\ \mathrm{t}}.
    \label{eq:M-diepgang}
\end{equation}
 
\subsection{Stabiliteitscontrole bij maximale diepgangsbelasting}
 
Bij $M_\mathrm{lading,T} = 38\,502$~kg:
 
\begin{align}
    T_\mathrm{gel}      &= 1{,}50 + \frac{38\,502}{77\,004} = 2{,}000\ \mathrm{m},
    \label{eq:Tgel-max}\\[4pt]
    \nabla_\mathrm{gel} &= \frac{111\,502}{1\,000} = 111{,}5\ \mathrm{m}^3,
    \label{eq:nabla-max}\\[4pt]
    KB_\mathrm{gel}     &= 2{,}000 \times 0{,}5798 = 1{,}160\ \mathrm{m},
    \label{eq:KB-max}\\[4pt]
    BM_\mathrm{gel}     &= \frac{101{,}1}{111{,}5} = 0{,}907\ \mathrm{m},
    \label{eq:BM-max}\\[4pt]
    KG_\mathrm{gel}     &= \frac{89\,790 + 42\,352}{111\,502} = 1{,}185\ \mathrm{m},
    \label{eq:KG-max}\\[4pt]
    GM_\mathrm{gel}     &= 1{,}160 + 0{,}907 - 1{,}185 = \mathbf{0{,}882\ \mathrm{m}}.
    \label{eq:GM-max}
\end{align}
 
De metacentrische hoogte bij maximale belasting ($GM = 0{,}882$~m) ligt ruimschoots
boven zowel het IMO-minimum (0{,}15~m) als de aanbevolen ondergrens (0{,}30~m).
De \textbf{bindende randvoorwaarde is de diepgangsgrens}, niet de stabiliteitseis.
 
\subsection{Overzicht resultaten}
 
\begin{table}[H]
    \centering
    \caption{Stabiliteitsoverzicht \emph{t~Jansje} – lege, geladen (ontwerpconfiguratie) en maximaal geladen conditie}
    \label{tab:stabiliteit-overzicht-max}
    \begin{tabular}{lrrrl}
        \toprule
        Grootheid & Leeg & Geladen (ontwerp) & Max.\ geladen & Eenheid \\
        \midrule
        Ladingmassa $M_\mathrm{lading}$          &  0{,}0  & 36{,}2  & 38{,}5  & t   \\
        Diepgang $T$                              &  1{,}500& 1{,}970 & 2{,}000 & m   \\
        Verdringingsvolume $\nabla$               & 73{,}0  &109{,}2  &111{,}5  & m$^3$ \\
        Totaalgewicht $W$                         & 73{,}0  &109{,}2  &111{,}5  & t   \\
        Hoogte drijfpunt $KB$                     &  0{,}870&  1{,}142&  1{,}160& m   \\
        Metacentrische straal $BM$                &  1{,}384&  0{,}926&  0{,}907& m   \\
        Metacentrum $KM = KB + BM$               &  2{,}254&  2{,}068&  2{,}067& m   \\
        Zwaartepunt $KG$                          &  1{,}230&  1{,}187&  1{,}185& m   \\
        Metacentrische hoogte $GM$                &  1{,}024&  0{,}881&  0{,}882& m   \\
        IMO-minimum $GM_\mathrm{min}$             & \multicolumn{3}{c}{0{,}150} & m \\
        Aanbevolen minimum $GM_\mathrm{aanbevolen}$& \multicolumn{3}{c}{0{,}300} & m \\
        Slingerperiode $T_\phi$                   &  3{,}37 &  3{,}63 &  –      & s   \\
        \bottomrule
    \end{tabular}
\end{table}
 
\section{Conclusie maximale belading}
\label{sec:max-belading-conclusie}
 
De maximaal toelaatbare extra ladingmassa voor \emph{t~Jansje} bedraagt
$M_{\mathrm{lading,max}} = \mathbf{38{,}5\ \mathrm{t}}$, begrensd door de
diepgangsgrens $T_\mathrm{max} = 2{,}00$~m. In dit scenario bedraagt
$GM_\mathrm{gel} = 0{,}882$~m, ruim boven het IMO-minimum van 0{,}15~m en de
aanbevolen binnenvaartondergrens van 0{,}30~m. De ontwerpconfiguratie van 1\,936
zeezoutbatterijcellen (36{,}2~t, $T_\mathrm{gel} = 1{,}97$~m, $GM = 0{,}881$~m) valt
daarmee veilig binnen de grens, met een resterende marge van
$38{,}5 - 36{,}2 = \mathbf{2{,}3\ \mathrm{t}}$.